A la fin de ce chapitre, je saurai :
\begin{competences}
    \point Établir l’équation d’onde dans le cas d’une corde infiniment souple dans l’approximation des petits mouvements transverses.
    \point Identifier une équation de d’Alembert.
    \point Exprimer la célérité en fonction des paramètres du milieu.
    \point Citer des exemples de solutions de l’équation de d’Alembert unidimensionnelle.
    \point Établir la relation de dispersion à partir de l’équation de d’Alembert.
    \point Utiliser la notation complexe.
    \point Définir le vecteur d’onde, la vitesse de phase.
    \point Décomposer une onde stationnaire en ondes progressives, une onde progressive en ondes stationnaires.
    \point Justifier et exploiter des conditions aux limites.
    \point Définir et décrire les modes propres.
    \point Construire une solution quelconque par superposition de modes propres.
    \point Associer mode propre et résonance en régime forcé.
    \point Décrire un câble coaxial par un modèle à constantes réparties sans perte.
    \point Établir les équations de propagation dans un câble coaxial sans pertes modélisé comme un milieu continu caractérisé par une inductance linéique et une capacité linéique.
    \point Établir l’expression de l’impédance caractéristique d’un câble coaxial.
    \point \faSignLanguage Étudier la réflexion en amplitude de tension pour une impédance terminale nulle, infinie ou résistive.
    \point Classer les ondes sonores par domaines fréquentiels.
    \point Justifier les hypothèses de l’approximation acoustique par des ordres de grandeur.
    \point Écrire les équations locales linéarisées : conservation de la masse, équation thermodynamique, équation de la dynamique.
    \point Établir l’équation de propagation de la surpression formulée avec l’opérateur laplacien.
    \point Exprimer la célérité en fonction de la température pour un gaz parfait.
    \point Citer les ordres de grandeur de la célérité pour l’air et pour l’eau.
    \point Utiliser les expressions admises du vecteur densité de courant énergétique et de la densité volumique d’énergie associés à la propagation de l’onde.
    \point Définir l’intensité sonore et le niveau d’intensité sonore.
    \point Citer quelques ordres de grandeur de niveaux d’intensité sonore.
    \point Décrire le caractère longitudinal de l'onde sonore.
    \point Discuter de la validité du modèle de l’onde plane en relation avec le phénomène de diffraction.
    \point Utiliser le principe de superposition des ondes planes progressives harmoniques.
    \point Établir et utiliser l’impédance acoustique définie comme le rapport de la surpression sur le débit volumique ou comme le rapport de la surpression sur la vitesse.
    \point Commenter l'expression fournie de la surpression générée par une sphère pulsante : atténuation géométrique, structure locale.
    \point \faSignLanguage Mettre en œuvre une détection synchrone pour mesurer une vitesse par décalage Doppler.
    \point Identifier les différents termes de l’équation locale de Poynting.
    \point Exprimer la puissance rayonnée à travers une surface à l’aide du vecteur de Poynting.
    \point Citer les domaines du spectre des ondes électromagnétiques et leur associer des applications.
    \point Établir les équations de propagation.
    \point Utiliser la notation complexe.
    \point Établir la relation entre le vecteur champ électrique, le vecteur champ magnétique et le vecteur d’onde.
    \point Associer la direction du vecteur de Poynting et la direction de propagation de l’onde.
    \point Associer le flux du vecteur de Poynting à un flux de photons en utilisant la relation d’Einstein-Planck.
    \point Citer quelques ordres de grandeur de flux énergétiques surfaciques moyens (laser hélium-néon, flux solaire).
    \point Utiliser le principe de superposition d’ondes planes progressives harmoniques.
    \point Identifier l'expression d'une onde électromagnétique plane progressive polarisée rectilignement.
    \point \faSignLanguage Utiliser des polariseurs et étudier quantitativement la loi de Malus.
\end{competences}